\chapter*{Abstract}
\label{ch:abstract}

Diese Arbeit untersucht den Einsatz von Transformer-Modellen zur automatisierten Erkennung von Gebärdensprache anhand von Landmark-Sequenzen, die mit MediaPipe extrahiert wurden. Ziel war die Entwicklung eines robusten Klassifikators, der mit Klassenimbalance und visuellen Ähnlichkeiten umgehen kann. Als Datengrundlage dient der PHOENIX-Datensatz, ergänzt durch abgeleitete Merkmale wie Geschwindigkeits- und Knochenvektoren. Die Methodik kombiniert gezieltes Feature-Engineering mit einer angepassten Transformer-Architektur und einer hierarchischen Inferenzstrategie. Die Experimente zeigen eine deutliche Steigerung der Genauigkeit gegenüber Standardarchitekturen und eine höhere Fehlertoleranz bei seltenen Klassen. Abschließend werden die wichtigsten Fehlerquellen analysiert und Ansätze zur Verbesserung diskutiert. Die Ergebnisse liefern Impulse für zukünftige Forschung und praktische Anwendungen in der Gebärdensprachübersetzung.